\documentclass[10pt,a4paper,openany,twoside, prof]{book}

\usepackage{layout_cours}

\begin{document}
\chapnumber{10} % numéro du chapitre
\titre{Vecteurs du plan et coordonnées} % titre du chapitre
\classe{Seconde} % classe
\entetegal


\section*{Corrections}
\subsubsection*{69 p.111}
$\v{AB} \coordvp{3-3}{-2-4} \iff \v{AB} \coordvp{0}{-6} $

$\v{BC} \coordvp{-7-3}{2-(-2)} \iff \v{BC} \coordvp{-10}{4} $

$\v{CA} \coordvp{3-(-7)}{4-2} \iff \v{CA} \coordvp{10}{2} $

On peut également appliquer Chasles, $\v{CA} = \v{CB} +\v{BA} = -\v{BC} - \v{AB}$

\subsubsection*{70 p.111}
$(\v{a}+\v{b}) \coordvp{3+2}{-4+3} \iff (\v{a}+\v{b}) \coordvp{5}{-1}$

$(\v{a}-\v{b}) \coordvp{3-2}{-4-3} \iff (\v{a}-\v{b}) \coordvp{1}{-7}$

$(2\v{a}) \coordvp{6}{8}$

$(-3\v{b}) \coordvp{-6}{-9}$

\subsubsection*{67 p.111}
$(\v{i}) \coordvp{1}{0}$

$(\v{j}) \coordvp{0}{1}$

$(\v{a}) \coordvp{1}{-3}$

$(\v{b}) \coordvp{2}{0}$

$(\v{c}) \coordvp{2}{3}$

$(\v{d}) \coordvp{-2}{3}$

$(\v{e}) \coordvp{-4}{1}$

$(\v{AB}) \coordvp{-3}{-1}$

$(\v{OC}) \coordvp{4}{2}$

$(\v{EF}) \coordvp{4}{-1}$

$(\v{KL}) \coordvp{0}{4}$

\subsubsection*{72 p.111}
\begin{enumerate}[label=(\roman*)]
\item
Soit $x,y \in\R$ les coordonnées de $A'$. C'est-à-dire $A'\coordpp{x}{y}$.

Alors $\vec{AA'}\coordvp{x-3}{y+2}$

On résout donc 

\begin{align*}
    \vec{AA'} =\vec{a} \iff \syst{x-3 &=& 3 \\ y+2 &=& 4} \iff \syst{x &=& 6 \\ y &=& 2}
\end{align*}

Donc $A'\coordpp{6}{2}$

\item
De même, soit $u,v \in\R$ les coordonnées de $B'$. C'est-à-dire $B'\coordpp{u}{v}$.

Alors $\vec{BB'}\coordvp{u+2}{v-1}$

On résout donc 

\begin{align*}
    \vec{BB'} =\vec{b} \iff \syst{u+2 &=& -1 \\ v-1 &=& 2} \iff \syst{u &=& -3 \\ v &=& 3}
\end{align*}

Donc $B'\coordpp{-3}{3}$

\end{enumerate}

\begin{rem}{}{}
    La lecture graphique n'est pas une méthode générale. Cela permet en revanche de vérifier son résultat.
\end{rem}

\subsubsection*{72 p.111}
\begin{enumerate}[label=(\alph*)]
    \item $\vec{EF}\coordvp{-5}{5}$
    \item $\vecu$ et $\vec{EF}$ sont colinéaires si et seulement s'il existe un coefficient de proportionnalité $k\in\R$ tel que $\vecu=k\vec{EF}$.
    
    $k=\frac{-3}{5}$ vérifie  $\vecu=k\vec{EF}=$.

    Donc les vecteurs sont colinéaires.

    \item 
    $\norm{\vec{EF}} = \sqrt{25+25} = 5\sqrt{2}$
    
    $\norm{\vecu} = \sqrt{9+9} = 3\sqrt{2}$
\end{enumerate}

\begin{rem}{}{}
    On peut aussi calculer la norme du vecteur $\vecu$ à partir de la norme du vecteur $\vec{EF}$.

    $\norm{\vecu} = \norm{k\vec{EF}} = |k|\times\norm{\vec{EF}} = \frac{3}{5} 5\sqrt{2} = 3\sqrt{2}$

    
\end{rem}
\subsubsection*{73 p.111}
Soit $x,y \in\R$ les coordonnées de $D$. C'est-à-dire $D\coordpp{x}{y}$.

Alors $\vec{DC}\coordvp{1-x}{3-y}$. D'autre part, $\vec{AB}\coordvp{4}{2}$

ABCD est un parallélogramme si et seulement si $\vec{AB}=\vec{DC}$

On résout donc :
\begin{align*}
    \vec{AB}=\vec{DC} \iff \syst{1-x &=& 4 \\ 3-y &=& 2} \iff \syst{x &=& -3 \\ y &=& 1}
\end{align*}

Donc $D\coordpp{-3}{1}$

\end{document}